Put \(N=M+1\). By \texttt{lem:finite-rank-staircase}, \(k=M+1=N\) is equivalent to \(\rho_{M,\alpha}=0\), and, because \(0<\alpha<1\), this is equivalent to \(0<\alpha<2/(M+1)\). Thus this is precisely the \(j=0\) step of the finite-rank staircase. By the convention in the corrected \texttt{def:rank-inversion},
\[r_{0,(k)}=r_{1,(k)}=+\infty.\]
Consequently the un-clipped residual-centered interval is all of \(\mathbb R\) in each arm. The reported interval is its intersection with the known outcome range, and hence
\[I_{a,w}=\mathbb R\cap[-K,K]=[-K,K],\qquad a\in\{0,1\}.\]
Taking the Minkowski difference gives
\[\mathcal C_w=[-K,K]-[-K,K]=[-2K,2K].\]
By \texttt{ass:bounded-potential-outcomes}, both \(Y_{0T}(1)\) and \(Y_{0T}(0)\) lie in \([-K,K]\) almost surely. Therefore
\[-2K\le Y_{0T}(1)-Y_{0T}(0)=\tau_T\le2K,\]
so \(\tau_T\in\mathcal C_w\) surely for every law and every pair of trained predictor maps. Conditioning on \(\mathcal T\) preserves this probability-one assertion.